\chapter{总结与展望}

\section{全文工作总结}
本文围绕流程工业场景下的领域自适应故障诊断问题，面向 Tennessee Eastman Process Domain Adaptation 数据集，初步完成了论文结构搭建、任务定义梳理与 benchmark 实验框架实现。通过统一数据接口、训练协议和结果分析流程，本文为后续比较不同领域自适应方法在跨工况故障诊断中的表现提供了基础支撑。

在方法层面，本文基于统一的一维卷积骨干网络，实现了 Source Only、CORAL、DAN、DANN、CDAN 与 DeepJDOT 等多种代表性方法，并围绕数据准备、模型训练、结果汇总和图表导出形成了较完整的实验工作流。当前阶段的主要成果在于完成了研究工作的整体组织，使论文写作、代码实现与后续实验分析能够围绕同一套结构持续推进。

\section{不足与未来展望}
当前工作仍存在若干不足。首先，论文中的大规模对比实验和定量结果分析尚需进一步补充，尤其需要在更多迁移场景下系统比较不同方法的稳定性与收益。其次，现阶段的研究重点仍以 benchmark 复现为主，对流程工业数据中时序依赖、类别不平衡和复杂工况变化的针对性建模还不够深入。最后，当前框架主要面向单源域到单目标域设置，尚未扩展到多源迁移、少样本目标域适配或更贴近工程部署的在线诊断情形。

后续研究可以从以下几个方向继续展开。其一，补充多场景实验结果，形成更完整的横向比较和误差分析。其二，在现有基线基础上探索更适合工业时序数据的结构设计与损失约束，例如结合时序上下文建模、类别级对齐或不确定性估计。其三，进一步完善实验平台的可扩展性，把多源域设置、消融实验和自动化报告能力纳入同一工作流。通过这些工作，可使本文从“完成 benchmark 搭建与方法复现”逐步推进到“形成具有针对性的流程工业领域自适应诊断研究结果”。
